\tocsection{\textbf{11/4/26 Lecture 28:} Zorn's Lemma and Topological Spaces}

\subsection*{Preliminaries from Logic}
\sectionspace

\begin{definition}[\textbf{Definition 28.1:} Cartesian Product]
    Let $X_{\alpha}$, $\alpha \in I$ be a collection of sets. The \emph{Cartesian product} of the $X_{\alpha}$ is
    \[ \prod_{\alpha \in I} X_{\alpha} = \left\{f : I \to \bigsqcup_{\alpha} X_{\alpha} \col f(\alpha) \in X_{\alpha}\right\} \]
\end{definition}

\textbf{Question:} If every $X_{\alpha}$ is nonempty, can $\prod_{\alpha \in I} X_{\alpha}$ be empty?\\

\textbf{Answer:} Whether or not it is, is independent of ZF (you need AC).

\begin{theorem}[\textbf{Axiom 28.2:} Axiom of Choice]
    Given any collection of nonempty sets $X_{\alpha}$, $\alpha \in I$, $\prod_{\alpha \in I} X_{\alpha}$ is nonempty.
\end{theorem}

\begin{definition}[\textbf{Definition 28.3:} Partially Ordered Sets]
    A \emph{partially ordered set (poset)} $(X, \preceq)$ is a set $X$ with a binary relation $\preceq$ that is reflexive, antisymmetric, and transitive. A poset is \emph{totally ordered} if for every $x, y \in X$, $x \preceq y$ or $y \preceq x$.
\end{definition}

\begin{example}{\textbf{Example:} Power Set}
    For any $X$ with at least two elements, $(\P(X), \subseteq)$ is a poset that is not totally ordered.
\end{example}

\begin{example}{\textbf{Example:} Trees}
    A tree $T$ can be made a poset by picking a ``root'' and saying that $x \preceq y$ if the unique path from the root to $y$ contains $x$.
\end{example}

\begin{definition}[\textbf{Definition 28.4:} Chains, Upper Bounds, and Maximal Elements]
    If $(X, \preceq)$ is a poset, a subset $Y \subseteq X$ is called a \emph{chain} if it is totally ordered. An \emph{upper bound} for $Y \subseteq X$ is an element $x \in X$ such that $y \preceq x$ for all $y \in Y$. An element $x \in X$ is \emph{maximal} if $x \preceq z \implies x = z$.
\end{definition}

\begin{lemma}[\textbf{Lemma 28.5:} Zorn's Lemma]
    Let $(X, \preceq)$ be a poset. Suppose every chain in $X$ has an upper bound. Then, $X$ has a maximal element.
\end{lemma}

\textbf{Fact:} Zorn's Lemma $\iff$ Axiom of Choice.

\begin{theorem}[\textbf{Theorem 28.6:} Basis Existance Theorem]
    Every vector space has a basis.
\end{theorem}
\begin{proof}{Proof:}
    Let $X = \{A \subseteq V \col A \tt{ linearly independent}\}$ with $\subseteq$. We need to show every chain has an upper bound. Let $Y$ be a chain in $X$. We claim
    \[ U = \bigcup_{B \in Y} B \]
    is an upper bound, which is to show it is linearly independent.\\


    Let $v_1, \ldots, v_r \in U$ be distinct. We want to show these are LI. Well, $v_i \in B_i \in Y$, so as $Y$ is totally ordered, all of them lie in some $B$. However, since $B$ is LI, $v_1, \ldots, v_r$ are LI.\\

    Thus, every chain has an upper bound, so by Zorn's Lemma, we have a maximal element $\tilde{A}$. Suppose $\tilde{A}$ is not a basis. Then, there exists $v \notin \Span(\tilde{A})$. Then, $\{v\} \cup \tilde{A}$ is LI, which contradicts that $\tilde{A}$ is maximal.

    $\qed$
\end{proof}

\begin{definition}[\textbf{Definition 28.7:} Topological Spaces]
    A \emph{topological space} $(X, \cal{T})$ is a set $X$ and a collection $\cal{T} \subseteq \P(X)$ called a \emph{topology} such that
    \begin{enumerate}[label = (\roman*)]
        \item $\emptyset, X \in \cal{T}$
        \item If $U_1, \ldots, U_n \in \cal{T}$, then $\bigcap_{i} U_i \in \cal{T}$
        \item If $U_{\alpha} \in \cal{T}$, then $\bigcup_{\alpha} U_{\alpha} \in \cal{T}$
    \end{enumerate}
    Elements of $\cal{T}$ are called \emph{open sets}.
\end{definition}

\begin{example}{\textbf{Example:} Discrete and Trivial Topologies}
    If $X$ is any set, $(X, \P(X))$ and $(X, \{\emptyset, X\})$ are topological spaces.
\end{example}

\begin{example}{\textbf{Example:} Metric Spaces}
    Given a metric space $(X, d)$, let $B_r(x) = \{y \col d(x, y) < r\}$. Then, say $U$ is open if for all $x \in U$, there exists $r > 0$ such that $x \in B_r(x) \subseteq U$.
\end{example}

\begin{definition}[\textbf{Definition 28.8:} Closed Sets]
    We say $A$ is \emph{closed} if $A^{c}$ is open.
\end{definition}

By De Morgan's laws, closed sets are preserved under finite unions and arbitrary intersections.

\begin{definition}[\textbf{Definition 28.9:} Closure, Interior, and Boundary]
    Given $A \subseteq X$, the \emph{closure} of $A$, $\overline{A}$, is the smallest closed set containing $A$. Similarly, the \emph{interior} of $A$, $A^{\circ}$, is the largest open set contained in $A$. The \emph{boundary} of $A$, $\partial A$, is $\overline{A} \setminus A^{\circ}$.
\end{definition}
